\documentclass[12pt]{article}
\usepackage{amsmath,amssymb,amsthm}
\usepackage[margin=1in]{geometry}
\usepackage{algorithm,algorithmic}

\title{Problem 266: Pseudo Square Root}
\author{Project Euler Solution}
\date{}

\begin{document}
\maketitle

\section{Problem Statement}
The divisors of 12 are: 1, 2, 3, 4, 6, 12. The largest divisor of 12 that does not exceed $\sqrt{12}$ is 3. We call the largest divisor of $n$ that does not exceed $\sqrt{n}$ the \emph{pseudo square root} (PSR) of $n$.

Let $P = \prod_{p < 190, p \text{ prime}} p$. Find $\text{PSR}(P)$. Give the last 16 digits.

\section{Mathematical Analysis}

\subsection{Setup}
The primes below 190 are:
\[ 2, 3, 5, 7, 11, 13, \ldots, 179, 181 \]
There are 42 such primes. Since $P$ is the product of distinct primes, it is squarefree, and every divisor of $P$ has the form $d = \prod_{p \in S} p$ for some subset $S$ of the primes.

We seek the largest such $d$ with $d \leq \sqrt{P}$, equivalently:
\[ \log d \leq \frac{1}{2}\sum_{p < 190}\log p \]

\subsection{Meet in the Middle}
With 42 primes, enumerating all $2^{42} \approx 4.4 \times 10^{12}$ subsets is infeasible. We split the primes into two halves $A$ and $B$ of 21 primes each, and use a \textbf{meet-in-the-middle} strategy:

\begin{enumerate}
    \item For group $A$, enumerate all $2^{21}$ subset products. Store $(\log d_A, \; d_A \bmod 10^{16})$.
    \item Sort these by $\log d_A$.
    \item For each of the $2^{21}$ subset products from group $B$, compute $\log d_B$ and $d_B \bmod 10^{16}$.
    \item Binary search group $A$ for the largest $\log d_A$ such that $\log d_A + \log d_B \leq \frac{1}{2}\log P$.
    \item Track the maximum $\log(d_A \cdot d_B)$ and the corresponding value modulo $10^{16}$.
\end{enumerate}

\subsection{Arithmetic Details}
Since $P$ is astronomically large, we use:
\begin{itemize}
    \item \textbf{Logarithms} for comparison: $\log(\sqrt{P}) = \frac{1}{2}\sum_p \log p$.
    \item \textbf{Modular arithmetic} for the answer: all products computed modulo $10^{16}$.
    \item \textbf{128-bit integers} (\texttt{\_\_int128}) for modular multiplication.
\end{itemize}

\section{Complexity}
\begin{itemize}
    \item Time: $O(2^{n/2} \cdot n/2)$ where $n = 42$, giving $\sim 44$ million operations.
    \item Space: $O(2^{n/2}) \approx 2$ million entries.
\end{itemize}

\section{Answer}

\[
\boxed{1096883702440585}
\]

\end{document}
